% Generated by experiments/make_numbers.py -- do not edit.
%
% Collections and named gene sets, from
% results/benchmark_{gse176078,gse161529}_{by_collection,by_theme}.csv.

\begin{table}[t]
\centering
\small
\begin{tabular}{lrrrrrr}
\toprule
Collection & Sets & Median $k$ & Detection & Median $|\rho|$ & Pairs $\geq 0.3$ & Within-compartment $\rho$ \\
\midrule
\multicolumn{7}{l}{\emph{GSE176078}} \\
Hallmark & 50 & 168 & 6.6\% & 0.349 & 57\% & +0.64 \\
Immunologic & 400 & 185 & 6.4\% & 0.273 & 45\% & +0.41 \\
Reactome & 250 & 26 & 7.0\% & 0.193 & 30\% & +0.43 \\
BioCarta & 288 & 14 & 8.4\% & 0.175 & 22\% & +0.30 \\
KEGG & 186 & 46 & 4.9\% & 0.173 & 27\% & +0.26 \\
WikiPathways & 200 & 27 & 5.6\% & 0.162 & 22\% & +0.28 \\
\midrule
\multicolumn{7}{l}{\emph{GSE161529}} \\
Hallmark & 50 & 176 & 5.8\% & 0.496 & 77\% & +0.72 \\
Immunologic & 400 & 193 & 5.1\% & 0.470 & 75\% & +0.73 \\
Reactome & 250 & 23 & 5.2\% & 0.259 & 44\% & +0.57 \\
KEGG & 186 & 52 & 4.0\% & 0.255 & 43\% & +0.42 \\
BioCarta & 292 & 15 & 6.6\% & 0.248 & 39\% & +0.48 \\
WikiPathways & 200 & 25 & 4.9\% & 0.245 & 42\% & +0.35 \\
\bottomrule
\end{tabular}
\caption{Every collection the benchmark draws on, scored in all
sixteen annotated compartments of GSE176078 and all seventeen of
GSE161529, so a difference between rows within a block is a difference
between sets rather than between the compartments that happen to carry
them. \emph{Collection} is the MSigDB collection the set names belong
to, with the immunologic signatures grouped by elimination; \emph{Sets}
counts the sets retained after the usable-overlap filter, \emph{Detection}
is the median fraction of cells in which a set's genes are seen (per
cent), \emph{Median $|\rho|$} the median absolute correlation of the
score with the cell's detected-gene count over that collection's
(compartment, set) pairs, and \emph{Within-compartment $\rho$} the
median over compartments of the Spearman correlation between a set's
median detection rate and its absolute score--depth correlation. That
last column is the diagnostic's within-compartment performance, and it
has the same sign in every collection, every compartment and every
cohort.}
\label{tab:generality}
\end{table}

\begin{table}[t]
\centering
\footnotesize
\setlength{\tabcolsep}{4pt}
\begin{tabular}{llrrrr}
\toprule
Theme & Gene set & $k$ & Detection & Median $|\rho|$, 1st & Median $|\rho|$, 2nd \\
\midrule
Hypoxia & \multicolumn{1}{p{6.1cm}}{\raggedright HALLMARK\_\allowbreak{}HYPOXIA} & 189 & 5.4\% & 0.466 & 0.649 \\
 & \multicolumn{1}{p{6.1cm}}{\raggedright GSE26023\_\allowbreak{}PHD3\_\allowbreak{}KO\_\allowbreak{}VS\_\allowbreak{}WT\_\allowbreak{}NEUTROPHIL\_\allowbreak{}HYPOXIA\_\allowbreak{}DN} & 178 & 3.7\% & 0.318 & 0.515 \\
 & \multicolumn{1}{p{6.1cm}}{\raggedright BIOCARTA\_\allowbreak{}P53HYPOXIA\_\allowbreak{}PATHWAY} & 21 & 11.8\% & 0.134 & 0.420 \\
Cell cycle & \multicolumn{1}{p{6.1cm}}{\raggedright REACTOME\_\allowbreak{}THE\_\allowbreak{}ROLE\_\allowbreak{}OF\_\allowbreak{}GTSE1\_\allowbreak{}IN\_\allowbreak{}G2\_\allowbreak{}M\_\allowbreak{}PROGRESSION\_\allowbreak{}AFTER\_\allowbreak{}G2\_\allowbreak{}CHECKPOINT} & 76 & 26.5\% & 0.504 & --- \\
 & \multicolumn{1}{p{6.1cm}}{\raggedright REACTOME\_\allowbreak{}G2\_\allowbreak{}M\_\allowbreak{}CHECKPOINTS} & 127 & 10.3\% & 0.452 & 0.713 \\
 & \multicolumn{1}{p{6.1cm}}{\raggedright KEGG\_\allowbreak{}CELL\_\allowbreak{}CYCLE} & 124 & 6.5\% & 0.302 & 0.558 \\
 & \multicolumn{1}{p{6.1cm}}{\raggedright REACTOME\_\allowbreak{}TP53\_\allowbreak{}REGULATES\_\allowbreak{}TRANSCRIPTION\_\allowbreak{}OF\_\allowbreak{}ADDITIONAL\_\allowbreak{}CELL\_\allowbreak{}CYCLE\_\allowbreak{}GENES\_\allowbreak{}WHOSE\_\allowbreak{}EXACT\_\allowbreak{}ROLE\_\allowbreak{}IN\_\allowbreak{}THE\_\allowbreak{}P53\_\allowbreak{}PATHWAY\_\allowbreak{}REMAIN\_\allowbreak{}UNCERTAIN} & 20 & 11.8\% & 0.203 & 0.250 \\
 & \multicolumn{1}{p{6.1cm}}{\raggedright HALLMARK\_\allowbreak{}G2M\_\allowbreak{}CHECKPOINT} & 190 & 7.3\% & 0.200 & 0.569 \\
 & \multicolumn{1}{p{6.1cm}}{\raggedright HALLMARK\_\allowbreak{}E2F\_\allowbreak{}TARGETS} & 195 & 6.7\% & 0.159 & 0.491 \\
 & \multicolumn{1}{p{6.1cm}}{\raggedright REACTOME\_\allowbreak{}TRANSCRIPTIONAL\_\allowbreak{}REGULATION\_\allowbreak{}BY\_\allowbreak{}E2F6} & 34 & 8.3\% & 0.112 & --- \\
 & \multicolumn{1}{p{6.1cm}}{\raggedright REACTOME\_\allowbreak{}PTK6\_\allowbreak{}REGULATES\_\allowbreak{}CELL\_\allowbreak{}CYCLE} & 6 & 8.9\% & 0.096 & 0.218 \\
 & \multicolumn{1}{p{6.1cm}}{\raggedright BIOCARTA\_\allowbreak{}SKP2E2F\_\allowbreak{}PATHWAY} & 10 & 4.9\% & 0.085 & 0.251 \\
 & \multicolumn{1}{p{6.1cm}}{\raggedright WP\_\allowbreak{}MIR124\_\allowbreak{}PREDICTED\_\allowbreak{}INTERACTIONS\_\allowbreak{}WITH\_\allowbreak{}CELL\_\allowbreak{}CYCLE\_\allowbreak{}AND\_\allowbreak{}DIFFERENTIATION} & 6 & 8.1\% & 0.064 & --- \\
Immune checkpoint & \multicolumn{1}{p{6.1cm}}{\raggedright REACTOME\_\allowbreak{}PD\_\allowbreak{}1\_\allowbreak{}SIGNALING} & 26 & 8.4\% & 0.262 & 0.205 \\
 & \multicolumn{1}{p{6.1cm}}{\raggedright BIOCARTA\_\allowbreak{}CTLA4\_\allowbreak{}PATHWAY} & 20 & 5.6\% & 0.217 & 0.178 \\
Interferon & \multicolumn{1}{p{6.1cm}}{\raggedright HALLMARK\_\allowbreak{}INTERFERON\_\allowbreak{}GAMMA\_\allowbreak{}RESPONSE} & 196 & 10.8\% & 0.490 & 0.526 \\
 & \multicolumn{1}{p{6.1cm}}{\raggedright WP\_\allowbreak{}TYPE\_\allowbreak{}II\_\allowbreak{}INTERFERON\_\allowbreak{}SIGNALING} & 34 & 14.7\% & 0.442 & --- \\
 & \multicolumn{1}{p{6.1cm}}{\raggedright HALLMARK\_\allowbreak{}INTERFERON\_\allowbreak{}ALPHA\_\allowbreak{}RESPONSE} & 95 & 12.4\% & 0.355 & 0.425 \\
 & \multicolumn{1}{p{6.1cm}}{\raggedright WP\_\allowbreak{}TYPE\_\allowbreak{}III\_\allowbreak{}INTERFERON\_\allowbreak{}SIGNALING} & 10 & 8.3\% & 0.184 & 0.199 \\
\bottomrule
\end{tabular}
\caption{The named sets in the four themes a reader is most likely to
test the framework on, selected from the collections by name rather than
chosen: every set whose name carries a hypoxia, cell-cycle, immune-
checkpoint or interferon term is here. \emph{Median $|\rho|$} is the
median absolute correlation of the score with the cell's detected-gene
count over the compartments of each cohort: the first is GSE176078, the
second GSE161529, whose cells detect fewer genes against a higher
ceiling (Table~\ref{tab:cohorts}), so the same arithmetic predicts
more damage there. A dash is a set that one cohort did not retain.}
\label{tab:themes}
\end{table}
